\fbox{clarify notion dimension of a set} because we need it for doing an identitification

In the notation of $Pfeffer$ we have
$A = D_J \cap cl (B_r(0))$, $A^- = cl (A)$, 
$A^\circ = int (A)$ et $A^\dot = Bd(A)$. 



- Fact 1: $(\Rn \setminus D_J) \cap (B_r(0))$ is thin (see lemma ...)

- Fact 2a): $cl ( D_J \cap cl(B_r(0))) \ni y \mapsto v e^{...}$  is a bounded vector field.
%
Idea of the proof: Take the norm and bound the exponential from above by 1, and from below by 0.

- Fact 2b): $cl ( D_J \cap cl(B_r(0))) \ni y \mapsto y e^{...}$  is a bounded vector field.
Idea of the proof: Take the norm, use the triangle inequality, and bound the two resulting term separately.

- Fact 3a):  $cl ( D_J \cap cl(B_r(0))) \ni v \mapsto y e^{...}$ is continuous. Therefore it is also continuous on $cl ( D_J \cap cl(B_r(0))) \cap S$ where $S$ is any slight set.
- Fact 3b):  $cl ( D_J \cap cl(B_r(0))) \ni y \mapsto y e^{...}$ is continuous. Therefore it is also continuous on $cl ( D_J \cap cl(B_r(0))) \cap S$ where $S$ is any slight set.
IDEA of the proof: INSIST that we are continuous on $\Rn$, therefore on $cl ( D_J \cap cl(B_r(0)))$ therefore on any other subset.

Fact 4): [Differentiability] the two vectors are differentiable on $int (D_J \cup cl(B_r(0))) \setminus T$ where $T$ is a thin subset. We choose $T = (\Rn \setminus D_J) \cap cl (B_r(0))$. The proof that this chosen  $T$ is  given by lemma...

Fact 5) $\emptyset$ is a slight subset of $D_J \cap cl (B_r(0))$. TO CHECK

Fact 6) $D_J \cap cl (B_r(0))$ is admissible. TO CHECK.

Lemma: The Lebesgue measure in dimension (n-1) of $(\Rn\setminus D_J) \cap B_r(0)$ is sigma-finite.
In other words, the set $(\Rn \setminus D_J) \cap (B_r(0))$ is thin.
PROOF: 
Step 1) Let $\Omega = B_r(0)$. Note that $\Omega$ is bounded, convex and open.  Define $u : \Omega \to \R$ by by $u(y) = J(y)$.
As $J: \Rn \to \R$ is convex, we can invoke (\cite{rockafellar1970convex}, Theorem 10.4) to conclude that its restriction on $\Omega$ is convex and Lipschitz continuous. Therefore $u$ is convex and Lipschitz continuous, i.e., there exist $L>0$ such that for all $(x,y) \in \Omega \times \Omega$ we have $|u(\bx)-u(\by)| \leqslant L\normsq{\bx-by}$.
Step 2) %We invoke (\cite{alberti1992singularities}, Theorem 4.1) of Albertini et al. 
Note that $u$ is semi-convex since it is convex. Therefore $u$ is Lipschitz and semi-convex. We can invoke (\cite{alberti1992singularities}, Theorem 4.1) to conclude that the following set
$$
\{
y \in \Omega \,|\, dim(\partial u(y)) \geq 1 
\}
$$
is countably $H^{n-1}$-rectifiable where $H^{n-1}$ denotes the Haussdorff measure blabla... . 
\fbox{Clarify $dim(\partial u(y))$}. 
\fbox{IDENTIFICATION} 
$$
(R^n \setminus D_J) \cap B_r(0) =
\{
y \in \Omega \,|\, dim(\partial u(y)) \geq 1 
\}
$$
Therefore 
$$
H^{n-1}((R^n \setminus D_J) \cap B_r(0))
$$
is countably $H^{n-1}$-rectifiable. 
In $\Rn$ the Hausdorff $H^{n-1} $measure is proportional of the Lebesgue measure in $\R^{n-1}$[ref].
\fbox{Set is countably $H^{n-1}$-rectifiable implies the Lebesgue measure in dimension $n-1$ of this set is sigma-finite. }
%
Therefore the Lebesgue measure in $n-1$ of the above set is sigma-finite. 



Invoking th 4.14 (Pfeffer). $A = D_J \cap cl (B_r(0))$ is admissible (fact 6). 
$S = \emptyset$ (check $\emptyset$ is slight subset of $int(A)$, fact 5). $T = (R^n \setminus D_J) \cap cl (B_r(0))$ is a thin subset of $int(A)$ (fact 1).
The vector field  is bounded on $int(A)$ (fact 2a and 2b), is continuous on $int(A) \setminus S$ (fact 3a and 3b), differentiable on $int(A)\setminus T$ (fact 4). 

%%%%%%%%%%%%%%%%%%%%%%%

